% =====================================================================
%  A.14 -- Adenda etica de la tercera ronda de trabajo de campo
%  Proyecto SIMPA -- Equipo AHMRV -- ISR-401 -- UTEQ -- 2026-2027 PPA
%  Compilacion: pdflatex A14_Adenda_Tercera_Ronda.tex  (x2)
% =====================================================================
\documentclass[11pt,a4paper]{article}

\usepackage[utf8]{inputenc}
\usepackage[T1]{fontenc}
\usepackage[margin=2.4cm]{geometry}
\usepackage{array}
\usepackage{booktabs}
\usepackage{longtable}
\usepackage[table]{xcolor}
\usepackage{enumitem}
\usepackage{tcolorbox}
\usepackage{fancyhdr}
\usepackage{titlesec}
\usepackage{lastpage}
\usepackage[hidelinks]{hyperref}

\definecolor{uteqverde}{HTML}{1F6F3C}
\definecolor{grisclaro}{HTML}{F2F2F2}

\newcolumntype{L}[1]{>{\raggedright\arraybackslash}p{#1}}
\newcolumntype{C}[1]{>{\centering\arraybackslash}p{#1}}

\titleformat{\section}{\normalfont\large\bfseries\color{uteqverde}}{\thesection.}{0.5em}{}
\titlespacing*{\section}{0pt}{12pt}{5pt}
\titleformat{\subsection}{\normalfont\normalsize\bfseries}{\thesubsection}{0.5em}{}

\pagestyle{fancy}
\fancyhf{}
\renewcommand{\headrulewidth}{0.4pt}
\lhead{\footnotesize\color{uteqverde}A.14 --- Adenda ética de la tercera ronda}
\rhead{\footnotesize SIMPA --- Equipo AHMRV}
\cfoot{\footnotesize Página \thepage\ de \pageref{LastPage}}

\setlength{\parskip}{4pt}
\setlength{\parindent}{0pt}

\begin{document}

\begin{center}
{\footnotesize UNIVERSIDAD TÉCNICA ESTATAL DE QUEVEDO\\
Facultad de Ciencias de la Computación --- Carrera de Software}\\[10pt]
{\Large\bfseries\color{uteqverde} ADENDA ÉTICA}\\[3pt]
{\large Tercera ronda de trabajo de campo}\\[2pt]
{\footnotesize Documento A.14 --- extensión del protocolo aprobado A.01}
\end{center}

\vspace{6pt}

\section*{Identificación del documento}
\renewcommand{\arraystretch}{1.25}
\begin{longtable}{|L{5.0cm}|L{10.6cm}|}
\hline
\rowcolor{grisclaro}\multicolumn{2}{|l|}{\textbf{Datos de identificación}}\\
\hline
Documento & A.14 --- Adenda ética de la tercera ronda de trabajo de campo \\
\hline
Proyecto & Sistema Inteligente de Mantenimiento de Palma Africana (SIMPA) \\
\hline
Equipo & AHMRV --- Paralelo 4to ``A'' \\
\hline
Asignatura & Ingeniería de Requerimientos (ISR-401) \\
\hline
Docente responsable & Ing. Gleiston Cicerón Guerrero Ulloa, PhD \\
\hline
Analista líder & Villafuerte Rosero Allan Noé \\
\hline
Organización cliente & Palmicultora M (seudónimo) \\
\hline
Segunda organización fuente & Extractora R (seudónimo) \\
\hline
Categoría de riesgo & C --- riesgo mínimo operativo \\
\hline
Protocolo base & A.01 --- Protocolo de investigación \\
\hline
Adendas previas & A.13 --- Adenda ética de la segunda ronda (emitida el 3 de agosto de 2026) \\
\hline
Fecha de emisión & \rule{5.5cm}{0.4pt} \\
\hline
\end{longtable}

\begin{tcolorbox}[colback=red!4,colframe=red!55!black,boxrule=0.8pt,arc=2pt,left=6pt,right=6pt,top=5pt,bottom=5pt]
\textbf{Condición de vigencia.} Esta adenda debe quedar firmada, fechada y cargada en el repositorio
\texttt{08\_Etica/} \textbf{con fecha anterior a la primera entrevista de la tercera ronda}. La guía de la
Entrega 4 (2B), sección 10.1, gatekeeper G8, condiciona la admisión de la entrega al cumplimiento de esta
secuencia. Ninguna entrevista de la tercera ronda debe iniciarse antes de que este documento esté firmado y
versionado.
\end{tcolorbox}

\section{Objeto de la adenda}

Esta adenda amplía el protocolo aprobado (documento A.01) para cubrir la tercera y última ronda de trabajo
de campo del proyecto SIMPA. Declara la incorporación de una población de participantes distinta de la de
las rondas anteriores, las técnicas adicionales que se aplicarán, el tratamiento de los datos y las medidas
de privacidad reforzadas. Todas las disposiciones del protocolo base y de la adenda A.13 permanecen
vigentes; esta adenda no las sustituye: las extiende.

\section{Motivo de la tercera ronda}

\subsection{Vacíos detectados en la codificación temática}

La codificación temática abierta de las ocho entrevistas de las rondas primera y segunda produjo 68 códigos
distintos distribuidos en 16 categorías. La distribución es marcadamente asimétrica: cinco categorías
concentran la mayor parte de los fragmentos codificados, mientras que las categorías directamente
relacionadas con la interacción entre la persona usuaria y el sistema permanecen escasamente pobladas.

\renewcommand{\arraystretch}{1.2}
\begin{longtable}{|L{5.2cm}|C{2.4cm}|L{7.4cm}|}
\hline
\rowcolor{grisclaro}\textbf{Categoría} & \textbf{Fragmentos} & \textbf{Estado} \\
\hline
Calidad de la fruta & 27 & Suficientemente cubierta \\
\hline
Detección de afectaciones & 19 & Suficientemente cubierta \\
\hline
Tecnología & 15 & \textbf{Cubierta sólo desde la perspectiva operativa} \\
\hline
Registro operativo & 12 & Suficientemente cubierta \\
\hline
Polinización / Organización del trabajo & 11 c/u & Suficientemente cubiertas \\
\hline
Reportes & 5 & \textbf{Vacío --- indicadores para decisión} \\
\hline
Remuneración & 4 & \textbf{Vacío --- transparencia percibida} \\
\hline
Trazabilidad & 3 & \textbf{Vacío} \\
\hline
Protección de datos & 1 & \textbf{Vacío crítico} \\
\hline
Cumplimiento normativo & 1 & \textbf{Vacío crítico} \\
\hline
\end{longtable}

Los códigos \texttt{USUARIO\_NO\_EXPERTO}, \texttt{CURVA\_APRENDIZAJE},
\texttt{RECOMENDACION\_ACCIONABLE}, \texttt{DIAGNOSTICO\_POR\_IMAGEN} y \texttt{MINIMIZACION\_DATOS}
aparecen una sola vez cada uno. Son precisamente los códigos que sustentan el requisito no funcional
RNF-16 (explicabilidad del diagnóstico asistido por inteligencia artificial para personal sin formación
técnica) y el requisito funcional RF-15 (registro y visualización de recorridos del personal, condicionado
a consentimiento). Ambos requisitos se encuentran hoy respaldados por evidencia insuficiente.

\subsection{Qué busca cubrir esta ronda}

\begin{enumerate}[leftmargin=*,itemsep=2pt]
  \item Comprensibilidad de la interfaz y de las etiquetas del sistema por parte de personas que no
        participaron en su especificación.
  \item Explicabilidad percibida del diagnóstico asistido (RNF-16): qué necesita ver una persona para
        confiar o desconfiar de una recomendación automática.
  \item Aceptabilidad del registro georreferenciado de recorridos (RF-15) y sensibilidad hacia la
        protección de datos personales.
  \item Utilidad percibida de los reportes e indicadores (RF-19) para la toma de decisiones.
  \item Curva de aprendizaje estimada y condiciones de adopción o rechazo de la herramienta.
  \item Contraste de un subconjunto de requisitos ya elicitados frente a una audiencia externa al proceso.
\end{enumerate}

\section{Población de la tercera ronda}

\subsection{Entorno universitario}

La tercera ronda se ejecuta principalmente en las instalaciones de la Universidad Técnica Estatal de
Quevedo. Los participantes previstos son personas del entorno universitario vinculadas a áreas de
agronomía, producción vegetal, agricultura o campos afines, en dos grupos: personal docente, ingeniería y
profesionales por una parte, y estudiantado de carreras relacionadas por otra.

\subsection{Declaración expresa sobre niveles heterogéneos de conocimiento}

\begin{tcolorbox}[colback=uteqverde!5,colframe=uteqverde,boxrule=0.8pt,arc=2pt,left=6pt,right=6pt,top=5pt,bottom=5pt]
Se declara de forma expresa que los participantes de esta ronda \textbf{no necesariamente poseen
experiencia directa en plantaciones de palma africana} ni el conocimiento operativo específico de los
participantes de las rondas anteriores. \textbf{Esta heterogeneidad es deliberada y constituye la variable
de estratificación de la ronda}, no una limitación no controlada.

El propósito no es replicar la elicitación de dominio ya realizada, sino incorporar un estrato de contraste
con distintos niveles de conocimiento previo, que permita evaluar la comprensibilidad del sistema, la
explicabilidad de sus salidas y la robustez de los requisitos ya especificados ante una audiencia que no
participó en su construcción.
\end{tcolorbox}

\subsection{Perfiles previstos y asignación de códigos}

\begin{longtable}{|C{2.1cm}|L{5.0cm}|L{4.2cm}|L{3.6cm}|}
\hline
\rowcolor{grisclaro}\textbf{Código} & \textbf{Perfil previsto} & \textbf{Organización} & \textbf{Técnicas previstas} \\
\hline
ENTR-09 & Docente / ingeniería del área agropecuaria & UTEQ & Entrevista + walkthrough \\
\hline
ENTR-10 & Docente / ingeniería del área agropecuaria & UTEQ & Entrevista + walkthrough \\
\hline
ENTR-11 & Profesional del área agrícola & UTEQ / externo & Entrevista + walkthrough \\
\hline
ENTR-12 & Profesional o docente del área agrícola & UTEQ / externo & Entrevista \\
\hline
ENTR-13 & Estudiantado de carrera afín & UTEQ & Entrevista + walkthrough \\
\hline
ENTR-14 & Estudiantado de carrera afín & UTEQ & Entrevista + walkthrough \\
\hline
ENTR-15 & Estudiantado de carrera afín & UTEQ & Entrevista + walkthrough \\
\hline
ENTR-16 & Estudiantado de carrera afín & UTEQ & Entrevista \\
\hline
\end{longtable}

\begin{tcolorbox}[colback=grisclaro,colframe=black!35,boxrule=0.6pt,arc=2pt,left=6pt,right=6pt,top=4pt,bottom=4pt]
\textbf{Regla de clasificación.} La adscripción de cada participante al estrato \emph{técnico} o
\emph{no técnico} \textbf{no se deriva de su cargo ni de su condición de docente o estudiante}, sino de la
experiencia efectivamente manifestada en las preguntas de caracterización (sección 6 de los guiones). Un
estudiante con experiencia familiar en producción agrícola puede quedar clasificado como técnico; un
docente de un área no agronómica puede quedar clasificado como no técnico. La clasificación se registra en
el acta de cada sesión junto con el criterio aplicado.
\end{tcolorbox}

\section{Técnicas de recolección}

\begin{enumerate}[leftmargin=*,itemsep=3pt]
  \item \textbf{Entrevista semiestructurada individual.} Duración prevista de 20 a 30 minutos. Dos
        variantes con un núcleo común de preguntas que permite la comparación entre estratos: Guion A
        (perfil profesional/docente) y Guion B (perfil estudiantil).
  \item \textbf{Walkthrough del prototipo funcional (MVP).} Sesión de observación de 10 a 15 minutos
        inmediatamente posterior a la entrevista, sobre dos o tres escenarios preseleccionados y trazados
        a la matriz de trazabilidad. El moderador interviene lo mínimo imprescindible: la explicación
        anticipada de cada pantalla invalidaría la observación de usabilidad.
  \item \textbf{Observación y notas de campo.} Registro escrito de dudas, vacilaciones y rutas de
        navegación seguidas por el participante.
\end{enumerate}

\subsection{Grabación de audio y video}

Se registra audio de la entrevista y, cuando el participante lo autorice de forma expresa, video de la
sesión. La autorización de grabación es independiente de la autorización de participación: un participante
puede aceptar la entrevista y negar la grabación en video, en cuyo caso la sesión se documenta únicamente
con audio y notas de campo. Durante el walkthrough se graba preferentemente la pantalla y la voz, no el
rostro.

\section{Tratamiento de los datos y seudonimización}

\begin{enumerate}[leftmargin=*,itemsep=2pt]
  \item Cada participante recibe un código \texttt{ENTR-XX} en el momento de la firma del consentimiento.
        A partir de ese instante, ningún artefacto del proyecto lo identifica por su nombre.
  \item La tabla que vincula el código con la identidad real es el único documento de reidentificación.
        Reside exclusivamente en la zona restringida [R], dentro del contenedor cifrado con AES-256, y no
        se versiona en texto plano en ningún punto del repositorio.
  \item Las transcripciones sustituyen nombres propios, topónimos precisos y razones sociales por
        seudónimos y códigos antes de su incorporación a la zona pública [P].
  \item Los datos que se depositen posteriormente en un repositorio abierto lo harán en su versión
        anonimizada, bajo licencia Creative Commons Atribución 4.0 Internacional, conforme a la
        autorización expresa recogida en el consentimiento A.03 v3.
  \item Se aplica el principio de minimización: no se solicita cédula, dirección, teléfono, correo
        electrónico ni ningún dato personal que no resulte indispensable para la finalidad declarada.
\end{enumerate}

\section{Medidas de privacidad reforzadas y medida correctiva sobre la zona pública}

\begin{tcolorbox}[colback=red!4,colframe=red!55!black,boxrule=0.8pt,arc=2pt,left=6pt,right=6pt,top=5pt,bottom=5pt]
\textbf{Medida correctiva declarada.} Una revisión interna del repositorio, posterior a la emisión de la
adenda A.13, detectó que en la zona pública subsistía material identificable: nombres manuscritos legibles
en copias de consentimientos, al menos una fotografía con rostro sin enmascarar y nombres propios
incorporados a nombres de archivo. En consecuencia, la afirmación contenida en la sección 2 de la adenda
A.13 según la cual ningún dato identificable se había publicado en la zona de acceso público
\textbf{no se corresponde con el estado verificado del repositorio en esa fecha}.

El equipo declara esta discrepancia de forma voluntaria, ejecuta el saneamiento de la zona pública con
carácter previo al inicio de la tercera ronda y adopta las medidas de captura que se enumeran a
continuación para impedir su repetición. Se deja constancia expresa de que la corrección se realiza por
iniciativa del equipo y no como respuesta a un requerimiento.
\end{tcolorbox}

Medidas aplicables desde el momento de la captura:

\begin{enumerate}[leftmargin=*,itemsep=2pt]
  \item Ningún archivo destinado a la zona pública incorpora el nombre del participante en su
        denominación. La nomenclatura obligatoria es \texttt{AAAA-MM-DD\_Tipo\_ENTR-XX\_Tecnica.ext}.
  \item No se publican rostros. Las fotografías de la zona pública se enmascaran antes de su carga, y se
        publican únicamente cuando documentan el entorno o la sesión sin permitir identificación.
  \item No se publican firmas, cédulas ni ningún número de documento. La copia pública del consentimiento
        lleva enmascarados el nombre, la firma y el documento de identidad, y muestra visible únicamente
        el código de participante.
  \item Se eliminan los metadatos EXIF, incluida la geolocalización, de toda imagen antes de su
        incorporación a la zona pública.
  \item No se publica ninguna combinación de atributos que permita la identificación indirecta
        (por ejemplo, cargo específico más unidad académica concreta más año de graduación).
  \item Los originales sin enmascarar, así como cualquier contenedor cifrado del proyecto ---incluidos los
        alojados en el repositorio audiovisual complementario--- permanecen exclusivamente en zona
        restringida o cifrada. La contraseña correspondiente \textbf{no se publica en ningún repositorio,
        público ni complementario}, y se entrega únicamente al docente responsable por el canal del
        Sistema de Gestión Académica.
\end{enumerate}

\subsection{Saneamiento del historial de control de versiones}

Con posterioridad a la corrección declarada en esta sección, una auditoría del repositorio determinó que
el historial de Git ---no la zona pública actual, sino versiones anteriores conservadas en el historial---
seguía reteniendo copias sin anonimizar de los consentimientos y fotografías corregidos: nombres reales
visibles en el nombre de archivo y en el contenido de al menos tres participantes de las rondas primera y
segunda.

El equipo reescribió dicho historial mediante la herramienta \texttt{git filter-repo}, retirando las rutas
que contenían las versiones sin anonimizar, y verificó la corrección mediante inspección directa de los
objetos del historial. Se detectó posteriormente que GitHub conserva de forma permanente e inmutable la
referencia de un Pull Request abierto sobre la rama afectada, la cual seguía sirviendo el contenido
original y no puede eliminarse mediante operaciones estándar de Git. Por esta limitación, ajena al control
del equipo, el proyecto se migró a un repositorio nuevo sin historial de Pull Requests:
\url{https://github.com/AlanNVR/SIMPA_ISR401}. El repositorio anterior quedó archivado en modo privado.

Esta operación queda registrada en \texttt{CHANGELOG.md} y en \texttt{08\_Etica/README\_Etica.md}, conforme
al principio de no ocultar la corrección: se documenta hacia adelante, sin modificar la adenda A.13 que
declaró el hallazgo original. El equipo considera esta desviación \textbf{subsanada y cerrada} a partir de
la fecha de emisión de la presente adenda. Esta declaración no constituye una nueva solicitud de
aprobación: documenta, con fines de trazabilidad, un hallazgo y su corrección ya ejecutada.

\section{Repositorio audiovisual complementario}

Debido al agotamiento de la cuota de almacenamiento de Git LFS del repositorio principal, los contenedores
audiovisuales de mayor tamaño se alojan en un repositorio complementario, conforme a la orientación
comunicada personalmente por el docente responsable. El repositorio principal conserva toda la evidencia
de control y trazabilidad: \texttt{fichas\_tecnicas.csv}, \texttt{checksums.sha256}, las transcripciones y
las copias públicas enmascaradas. El procedimiento de acceso, descifrado y verificación de integridad se
documenta en \texttt{02\_Evidencias/00\_Restringido/README\_Evidencias\_Externas.md}.

\section{Riesgos identificados y mitigaciones}

\renewcommand{\arraystretch}{1.2}
\begin{longtable}{|C{1.5cm}|L{5.6cm}|L{8.0cm}|}
\hline
\rowcolor{grisclaro}\textbf{ID} & \textbf{Riesgo} & \textbf{Mitigación} \\
\hline
R-14.1 & Identificación indirecta de un participante por la combinación de perfil y unidad académica &
Los perfiles se publican en categorías amplias (docente/profesional del área agropecuaria; estudiantado de
carrera afín). No se publica carrera concreta, semestre ni asignatura. \\
\hline
R-14.2 & Incomodidad del participante ante la grabación &
Autorización de grabación separada y revocable. Posibilidad de participar con audio únicamente o sin
grabación. Se informa antes de iniciar. \\
\hline
R-14.3 & Percepción de evaluación o examen por parte del estudiantado &
Se declara al inicio de cada sesión que no existe respuesta correcta, que no se evalúa a la persona
participante y que la participación no tiene relación alguna con su calificación académica. \\
\hline
R-14.4 & Relación de asimetría entre estudiantado entrevistado y equipo entrevistador &
Ningún integrante del equipo entrevista a estudiantes sobre los que tenga cualquier función de evaluación.
La participación no se solicita durante horas de clase ni a través de docentes con poder de calificación
sobre la persona. \\
\hline
R-14.5 & Fatiga de la sesión al encadenar entrevista y walkthrough &
Duración total acotada a 45 minutos. El walkthrough es opcional y se ofrece tras la entrevista, no antes.
El participante puede detenerlo en cualquier momento. \\
\hline
R-14.6 & Atribución de credibilidad a un diagnóstico automático simulado &
Véase la sección 9. Se declara la simulación al cierre de cada sesión de walkthrough. \\
\hline
R-14.7 & Contaminación del análisis de saturación por mezcla de poblaciones &
Véase la sección 10. Los estratos se analizan por separado. \\
\hline
R-14.8 & Reaparición de datos identificables en la zona pública &
Checklist de cierre por participante y verificación de nomenclatura antes de cada carga al repositorio. \\
\hline
\end{longtable}

\section{Declaración sobre las funcionalidades simuladas del prototipo}

\begin{tcolorbox}[colback=red!4,colframe=red!55!black,boxrule=0.8pt,arc=2pt,left=6pt,right=6pt,top=5pt,bottom=5pt]
Dos funcionalidades del prototipo funcional presentan interfaz completa con resultado no real: la detección
de plagas por análisis de imagen (RF-07) y el conteo georreferenciado de flores polinizadas (RF-14). En
ambos casos la salida se genera por una función de ejemplo y no por inferencia sobre datos reales.

Durante el walkthrough esta condición \textbf{no se revela con antelación}, porque el objeto de observación
es precisamente qué elementos de la interfaz sustentan o erosionan la confianza de la persona usuaria en una
salida automática. La condición de simulación \textbf{se revela de forma expresa al cierre de cada sesión},
antes de la despedida, y se deja constancia de la revelación en el acta correspondiente. Ningún dato
recogido sobre confianza en el diagnóstico se reportará sin declarar simultáneamente que la salida
observada era simulada.
\end{tcolorbox}

\section{Relación con las rondas anteriores y con el análisis de saturación}

La primera y la segunda ronda elicitaron requisitos de dominio con participantes de Palmicultora M y
Extractora R (ENTR-01 a ENTR-08). La tercera ronda no sustituye ni replica esa elicitación: incorpora un
estrato de contraste y validación (ENTR-09 a ENTR-16).

\begin{tcolorbox}[colback=grisclaro,colframe=black!35,boxrule=0.6pt,arc=2pt,left=6pt,right=6pt,top=4pt,bottom=4pt]
\textbf{Consecuencia metodológica que se declara por anticipado.} La curva de saturación de significado se
calculará y reportará \textbf{por estrato}, no sobre la mezcla indiferenciada de las dieciséis entrevistas.
Agregar en una sola curva dos poblaciones con distinta relación con el dominio produciría una inflexión
artificial: los participantes universitarios generan pocos códigos nuevos de dominio por construcción, no
por saturación alcanzada. Se reportará la curva del estrato de dominio (ENTR-01 a ENTR-08), la del estrato
de contraste (ENTR-09 a ENTR-16) y la agregada, con la interpretación de cada una. Esta decisión se
incorporará como amenaza a la validez de constructo en el manuscrito.
\end{tcolorbox}

Un subconjunto de preguntas del guion de la tercera ronda se mantiene deliberadamente común con las rondas
anteriores, para permitir el contraste de códigos ya existentes. El resto se orienta a las categorías con
vacío declaradas en la sección 2.

\section{Vigencia}

Esta adenda entra en vigor en la fecha de su firma y permanece vigente hasta el cierre del Proyecto Fin de
Curso. No modifica la categoría de riesgo del proyecto, que se mantiene en C --- riesgo mínimo operativo.
No se recogen categorías especiales de datos personales. Toda persona participante conserva el derecho de
retirar su consentimiento en cualquier momento, con la consiguiente eliminación de su material de todos los
artefactos del proyecto.

\vspace{18pt}

\begin{center}
\begin{tabular}{C{6.6cm} C{6.6cm}}
\rule{6.2cm}{0.4pt} & \rule{6.2cm}{0.4pt} \\[2pt]
\textbf{Villafuerte Rosero Allan Noé} & \textbf{Ing. Gleiston C. Guerrero Ulloa, PhD} \\
Analista líder --- Equipo AHMRV & Docente responsable --- ISR-401 \\[16pt]
\multicolumn{2}{c}{Fecha: \rule{4.5cm}{0.4pt}} \\
\end{tabular}
\end{center}

\end{document}
